% answers_short.tex —— 1986 年·全国卷·理科 · 简短答案速查
% 由 tools/make_short_answers.py 自 answers.tex 的 \gksolhead/\gkanswer 提取，
% 并与 1986/1986_数学真题_全国理_答案要点.md 逐题比对后生成，勿手改（改源 MD 后重跑）。

\begin{gknotice}
\textbf{使用说明}：本节只给最终答案便于快速核对，完整推导见解析版。答案为编写者独立推导，\textbf{非官方标准答案}。
\end{gknotice}

\gkpart{选择题}
\begin{center}\begin{tabular}{ccccccccccc}
\toprule
题号 & 1 & 2 & 3 & 4 & 5 & 6 & 7 & 8 & 9 & 10 \\
\midrule
答案 & B & C & B & A & B & D & A & A & C & C \\
\bottomrule
\end{tabular}\end{center}

\gkpart{解答题}
\begin{itemize}[leftmargin=2.6em]
  \item[\textbf{11.}] $x=\dfrac12$ 或 $x=-\dfrac32$。
  \item[\textbf{12.}] $0$。
  \item[\textbf{13.}] $\dfrac{25\pi}{3}$。
  \item[\textbf{14.}] $\dfrac13$。
  \item[\textbf{15.}] $-40$。
  \item[\textbf{16.}] $\dfrac{11}{16}$。
  \item[\textbf{17.}] $\text{\text{平面} }PAC\perp\text{\text{平面} }PBC$。
  \item[\textbf{18.}] $(-\pi,-3)\cup(2\pi,7)$。
  \item[\textbf{19.}] 若 $OA=a$、$OB=b$，则 $C=(\sqrt{ab},0)$。
  \item[\textbf{20.}] 固定源字面前提矛盾，暂不确定。
  \item[\textbf{21.}] 见解析正文。
  \item[\textbf{22.}] 见解析正文。
\end{itemize}

\gkpart{附加题}
\begin{itemize}[leftmargin=2.6em]
  \item[\textbf{23.}] $y'=\arctan x^2+\dfrac{2x^2}{1+x^4}$。
  \item[\textbf{24.}] $y=x+1$。
\end{itemize}
